\Needspace{12\baselineskip}
Träte nun jedes Jahr \textit{ein} neues Mitglied ein, dessen Berechtigung im Durchschnitt zu dem obigen Geldwert angesetzt werden müsste, so ist klar, dass zur Deckung dieser immer neu geborenen Ansprüche die Zinsen eines eisernen Kapitals erforderlich und ausreichend wären, dessen Höhe
\begin{center}
\begin{tabular}{l}
bei dem Zinsfuß von \(3\tfrac{1}{2}\) Prozent zu 17182,86Talern\\
bei dem Zinsfuß von 4 Prozent zu 12912,50 Talern
\end{tabular}
\end{center}
\noindent angesetzt werden müsste.

Bis hierher ruht die Rechnung auf einer sicheren Grundlage. Zu ihrer Vollendung fehlt nun aber noch ein wesentliches Element, für das jede im Voraus gewagte Schätzung der Natur der Sache nach nur sehr unsichere Gültigkeit haben kann: Ich meine die durchschnittliche Zahl der künftig jedes Jahr neu beitretenden Teilnehmer. Befragt man die Erfahrung, so zeigt sich in neuerer Zeit eine sehr starke Vergrößerung dieses Elements. Für die ersten achtzig Jahre des Bestehens der Anstalt kann man die durchschnittliche Zahl der jährlich Beitretenden mit nahezu \(1\tfrac{1}{2}\) ansetzen; allein in den letzten 20 Jahren sind insgesamt 64 Beitritte\footnote{Dabei ist G. mitgezählt, dessen Witwe Pension bezieht, der selbst aber eigentlich nicht aufgenommen worden war. Muss man jedoch voraussetzen, dass die Verwaltung in einem ähnlichen künftigen Fall wieder einen ähnlichen Beschluss fassen wird, so dürfen solche Fälle, auch wenn sie selten sind, nicht unberücksichtigt bleiben, sondern müssen in einer längeren Reihe von Erfahrungen als \textit{auch} mögliche Fälle mitgezählt werden.} erfolgt, davon 23 im Jahrzehnt 1825—1835 und 41 im Jahrzehnt 1835—1845. Der Durchschnitt aus den letzten 20 Jahren wäre demnach \(3\tfrac{1}{5}\), oder, wenn man G. nicht mitzählen will, \(3\tfrac{3}{20}\). Für die letzten 10 Jahre allein beträgt dieser Durchschnitt \(4\tfrac{1}{10}\), und für die letzten 5 Jahre allein, in denen 27 neue Mitglieder eingetreten sind, sogar \(5\tfrac{2}{5}\). Daran erkennt man, wie misslich es ist, hieraus eine Vorhersage für die Zukunft abzuleiten. Ich würde jedoch nicht wagen, die künftige jährliche Durchschnittszahl niedriger als \(2\tfrac{1}{2}\) anzusetzen. Dieses Verhältnis könnte auch als durchaus stimmig mit der jetzigen Teilnehmerzahl 51 angesehen werden, insofern die bisherigen Erfahrungen 20 Jahre oder etwas mehr als durchschnittliche Dauer der Teilnahme ergeben haben. Unter dieser Annahme ergibt sich also der gegenwärtige Geldwert
\SourceLeaf{161}
der Verpflichtungen der Kasse gegenüber allen künftigen Teilnehmern, also der Verpflichtungen unter der dritten Rubrik:
\begin{center}
\begin{tabular}{l}
zu 42957 Talern bei \(3\tfrac{1}{2}\) Prozent Zinsfuß\\
zu 32281 Talern bei 4 Prozent Zinsfuß.
\end{tabular}
\end{center}
\noindent\hspace*{\parindent}Wollte man die jährlichen Beitritte aber im Durchschnitt mit 3 ansetzen, so müsste man für diese Rubrik je nach dem Zinsfuß 8591 oder 6456 Taler \textit{mehr} veranschlagen.

\Needspace{12\baselineskip}
Fassen wir schließlich die Ergebnisse der bisherigen Ansätze zusammen, so ergibt sich aus den drei Rubriken
\begin{center}
\footnotesize\setlength{\tabcolsep}{3pt}
\begin{tabular}{l r|r}
&\(3\tfrac{1}{2}\) Prozent&4 Prozent\\
\shortstack[l]{Verpflichtungen gegenüber den jetzigen\\Witwen . . . . . . . . . . . . . . . . .}&46529 Thl.&44582 Thl.\\
\shortstack[l]{Verpflichtungen gegenüber den Hinterbliebenen\\der jetzigen Teilnehmer}&60724 ,, &54014 ,, \\
\shortstack[l]{Verpflichtungen gegenüber den künftigen\\Teilnehmern . . . . . . . . . . . . .}&42957 ,, &32281 ,, \\\cline{2-3}
\shortstack[l]{Gesamtsumme aller Verpflichtungen\\ . . . . . . . . . . . . . . . . . . . . . . . .}&150210 Thl.&130877 Thl.
\end{tabular}
\end{center}
\noindent\hspace*{\parindent}Von den Nebenkosten will ich hier nur die folgenden Rubriken nach dem Durchschnitt der letzten 15 Jahre in Rechnung stellen (alles auf Gold umgerechnet):
\begin{center}
\small\setlength{\tabcolsep}{3pt}
\begin{tabular}{l r}
Bau- und Reparaturkosten, jährlich . . . . .&43 Rthl.\\
Prozesskosten nach Abzug der Erstattungen&39 ,, \\
Kosten der Rechnungsführung . . . . . . . . .&124 ,, \\\cline{2-2}
\multicolumn{1}{r}{Zusammen . . . . . .}&206 Rthl.
\end{tabular}
\end{center}
\noindent\hspace*{\parindent}Die Bau- und Reparaturkosten sind nach diesem Durchschnitt vielleicht zu niedrig angesetzt. In dem oben S. [132] angeführten P. M. werden sie für einen Zeitraum von 24 Jahren (1768—1792) mit 1844 Talern Kassenmünze (1975 Thl. 17 Ggr. Gold) angesetzt, was somit einen jährlichen Durchschnittsbetrag von 82 Talern ergeben würde: Ich will jedoch bei dem obigen Mittel bleiben. Diese 206 Taler jährlicher Ausgaben entsprechen einem Kapitalbedarf von 5886 oder 5150 Talern, je nach dem Zinsfuß.

Es darf nicht übersehen werden, dass das allen Rechnungen zugrunde liegende Interusurienprinzip nur insoweit gilt, als alle Zinsen rechtzeitig eingehen und sofort wieder verwendet werden. In der Praxis ist dies aber nicht ausführbar; vielmehr liegt immer ein bald kleinerer, bald größerer Teil des Vermögens ungenutzt. Den mittleren Betrag dieser Geldsumme könnte man der Kürze halber Betriebsfonds nennen, obwohl diese Bezeichnung genau genommen nur auf den Teil des baren Geldwerts passt, der bereitgehalten werden muss, damit die Kasse ihren eigenen Zahlungsverpflichtungen pünktlich nachkommen kann. Dazu muss aber auch das Geld gerechnet werden, das zeitweise ungenutzt liegt, wenn sich nicht sofort eine Gelegenheit zu sicherer und vorteilhafter Anlage bietet, sowie die ausstehenden Rückstände. Wie viel dies alles im tatsächlichen Durchschnitt ausmacht, lässt sich den Jahresrechnungen nur unvollkommen entnehmen, da diese nur angeben, \textit{was} im Laufe des Jahres eingenommen und ausgegeben wurde, aber nicht an welchem Datum (dass dies \textit{zum Teil} aus den \textit{Belegen}, und auch aus diesen nicht ohne Ungewissheit, ergänzt werden \textit{könnte}, kommt hier nicht in Betracht). Ich habe jedoch versucht, diesen Betrag aus den Rechnungsabschlüssen der letzten 11 Jahre\footnote{Früher stellten die Rechnungen den Vermögensstand überhaupt nicht auf.} so gut wie möglich abzuleiten, und danach 2636 Taler als mittleren Betrag des nicht zinstragenden Vermögensteils gefunden.

\Needspace{10\baselineskip}
Ich lasse nun die Zusammenfassung dieser verschiedenen Posten folgen:
\SourceLeaf{162}
\begin{center}
\fontsize{8}{9.6}\selectfont\setlength{\tabcolsep}{3pt}
\begin{tabular}{l r|r}
&\(3\tfrac{1}{2}\) Prozent&4 Prozent\\
\shortstack[l]{Verpflichtungen gegenüber Witwen, jetzigen\\und künftigen Teilnehmern}&150210 Thl.&130877 Thl.\\
\shortstack[l]{Fonds zur Bestreitung der Bau-, Prozess-\\und Rechnungskosten . .}&5886 ,, &5150 ,, \\
Unverzinslicher Fonds . . . . . . . . . . . . . . . . . . . . . . . . . . . . . .&2636 ,, &2636 ,, \\\cline{2-3}
\multicolumn{1}{r}{Gesamtsumme . . . . .}&158732 Thl.&138663 Thl.
\end{tabular}
\end{center}

